\section{Search Insert Position}
\label{sec:bs-search-insert}

\problemstatement{We are given a sorted array $\textit{arr}$ of $n$
distinct integers and a target value $X$. If $X$ is found in the array,
return its index. If not, return the index at which $X$ \textbf{would be
inserted} to keep the array sorted.}

\begin{patternbox}
This problem is worth pausing on specifically \emph{because} it looks new
but is not: \emph{``the index where $X$ would be inserted to keep the array
sorted''} is, word for word, a description of the \textbf{lower bound} of
$X$ from Section~\ref{sec:bs-lower-bound}. The position just before the
first element $\ge X$ is exactly where $X$ slots in. Recognizing when a
new-sounding problem statement is a known pattern in disguise is as
important a skill in this chapter as deriving a pattern from scratch.
\end{patternbox}

\subsection*{Understanding the problem carefully}

If $X$ is present in the array, the lower bound of $X$ is precisely the
index of $X$ (since the array has distinct elements, the first index with
$\textit{arr}[i] \ge X$ is the unique index where $\textit{arr}[i] = X$).
If $X$ is absent, the lower bound is the first index whose value exceeds
$X$ -- which is exactly the index where $X$ should be inserted, since
everything before that index is smaller than $X$ and everything from that
index onward is larger.

\begin{ideabox}[Candidate Space and Predicate]
Identical to lower bound: candidate space is indices $0, \dots, n$,
predicate $P(i) = (\textit{arr}[i] \ge X)$, answer is the smallest index
where $P$ holds.
\end{ideabox}

\subsection*{Why no new derivation is needed}

There is nothing new to prove here -- this section exists to train the
habit of \emph{reducing} an unfamiliar-sounding problem to a pattern you
already have a proven, working template for, rather than re-deriving binary
search bounds from scratch every time a problem uses different words for
the same underlying question.

\subsection*{Algorithm}

Identical to Algorithm~\ref{sec:bs-lower-bound}'s lower-bound procedure,
called directly on $X$.

\subsection*{Dry run}

\dryrun{Take $\textit{arr} = [1, 3, 5, 6]$, $X = 5$.}

This is present at index $2$; the lower bound of $5$ is index $2$
(verify: $\textit{arr}[2]=5 \ge 5$, and $\textit{arr}[1]=3 < 5$), so the
answer is $2$.

Now take $X = 2$ on the same array. $2$ is absent. The lower bound of $2$
is index $1$ (the first element $\ge 2$ is $\textit{arr}[1]=3$), so $2$
would be inserted at index $1$, giving $[1,2,3,5,6]$ -- correct.

\subsection*{C++ implementation}

\begin{lstlisting}
#include <bits/stdc++.h>
using namespace std;

int searchInsert(vector<int>& arr, int X) {
    int n = arr.size();
    int lo = 0, hi = n - 1, ans = n;

    while (lo <= hi) {
        int mid = lo + (hi - lo) / 2;
        if (arr[mid] >= X) {
            ans = mid;
            hi = mid - 1;
        } else {
            lo = mid + 1;
        }
    }
    return ans;
}

int main() {
    vector<int> arr = {1, 3, 5, 6};
    cout << "Insert position for 5: " << searchInsert(arr, 5) << endl;
    cout << "Insert position for 2: " << searchInsert(arr, 2) << endl;
    return 0;
}
\end{lstlisting}

\begin{complexitybox}
\textbf{Time Complexity.} $O(\log n)$.
\textbf{Space Complexity.} $O(1)$.
\end{complexitybox}

\begin{takeawaybox}
Before deriving a new binary search predicate from scratch, check whether
the problem is just lower bound, upper bound, or a simple combination of
the two, restated in different words. Many ``new'' binary search problems
on interview sheets are exactly this.
\end{takeawaybox}
